% ID: SEL-D0009-BC-E05
% ESTADO: cotejada. Segunda pasada visual completa y borrador compilado revisados.
% Fuente: Selectividad/matematicas_aplicadas/ORD_RESER_Modelo-3-Enunciados (Mat-Ap-2021).pdf
% SHA-256: 1dc0455361ce91a1157a7b3e4a06db95bc988e8dd0c9756e9aeb74b8cbb73a81
% Página: 2; bloque C; ejercicio original 5.

Sean $A$ y $B$ dos sucesos de un mismo experimento aleatorio de los que
se sabe que:
\[
P(A-B)=0.3\qquad P(A^c)=0.35\qquad P(B)=0.55
\]
\begin{enumerate}[label=\alph*)]
  \item \textbf{(0.8 puntos)} Calcule la probabilidad de que suceda al menos
  uno de ellos.
  \item \textbf{(0.6 puntos)} Calcule la probabilidad de que ocurra $B$,
  sabiendo que no ha ocurrido $A$.
  \item \textbf{(0.6 puntos)} Calcule la probabilidad de que no ocurra
  ninguno de los dos sucesos.
  \item \textbf{(0.5 puntos)} Razone si los sucesos $A$ y $B$ son independientes.
\end{enumerate}
